We now introduce the notion of a simplicial subset, along with two important constructions known as the \emph{horn}, and the \emph{boundary/simplicial sphere}, which are of particular importance in a variety of contexts.

\begin{definition}\label{3.4.1}
    Let $X_\bullet$ be a simplicial set. A \emph{simplicial subset} of $X_\bullet$ is a simplicial set $Y_\bullet$ where
    \begin{itemize}
        \item For all $n$, the $n$-simplices of $Y_\bullet$ are a subset of those of $X_\bullet$, meaning $Y_n\subseteq X_n$;
        \item The faces and degeneracies in $Y_\bullet$ agree with those in $X_\bullet$; in particular, the faces $d_i^Y$ of $Y_\bullet$ are defined as the restrictions of the faces $d_i^X$ and degeneracies $s_i^X$ of $X_\bullet$ along the subset $Y_n\subseteq X_n$
        \[
        d_i^Y = d_i^X|_{Y_n}, \qquad s_i^Y = s_i^X|_{Y_n}.
        \]
    \end{itemize}
    We write $\Y\subseteq\X$ to denote that it is a simplicial subset.
\end{definition}

\begin{remark}\label{3.4.2}
    Unfortunately, not all collections of subsets $Y_n\subset X_n$ of a simplicial set $X_\bullet$ form a simplicial subset. Yes, one can define the faces and degeneracies for $Y_\bullet$ by restricting those of $X_\bullet$ along the relevant subsets, but \emph{the image of each restriction must be contained in the corresponding subset $Y_n$}.
    
    For example, consider the discrete simplicial set $X_\bullet$ where $X_n=\mathbb{Z}$ for all $n$, and all faces and degeneracies are the identities. We can construct subsets $Y_\bullet$ by taking $Y_n=X_n$ for all $n$ \emph{except} for $n=1$, where we take $Y_n=\mathbb{N}$. It doesn't matter whether or not we choose $0\in\mathbb{N}$ here. By \cref{3.4.1}, for $Y_\bullet$ to be a simplicial subset, the faces and degeneracies must agree with those of $X_\bullet$; in particular they are the restrictions along the subsets $Y_\bullet$. Since $Y_0=X_0$, the face $d_0 : Y_0 \to Y_1$ must be the same. Since the face is the identity, it maps each element to itself. Note that $-1\in Y_0=\mathbb{Z}$, and $d_0(-1)=-1\notin Y_1=\mathbb{N}$. Since the image is not in the subset $Y_1$, $\Y$ cannot be a simplicial subset, since the definition of $\Y$ as given does not form a valid simplicial set.
\end{remark}

\begin{remark}\label{3.4.3}
    In contrast to \cref{3.4.2}, it turns out that all one has to check to determine whether an object $Y_\bullet$ is a simplicial subset of a simplicial set $\X$ is
    \begin{itemize}
        \item $Y_n\subseteq X_n$ for all $n$;
        \item The faces and degeneracies agree with $X$;
        \item The faces and degeneracies are well-defined in the sense of \cref{3.4.2}, meaning the sets $Y_n$ contain the images of all faces and degeneracies.
    \end{itemize}
    This is fairly easily shown. The first two criteria guarentee that if $\Y$ is a simplicial set, it will be a simplicial subset of $\X$, so we need only check that it is a simplicial set. Since the faces and degeneracies are all well-defined, \cref{3.1.7} tells us that we need only check the faces and degeneracies satisfy the simplicial identities. This is manifestly true: the faces and degeneracies in $\Y$ are simply restrictions of those on $\X$, and since $\X$ is a simplicial set, the faces and degeneracies satsify the simplicial identities on the whole of $\X$. This means they must satisfy them for every element of each $X_n$, and thus every subset of each $X_n$.

    In particular, we see that to specify a simplicial subset $\Y$ of $\X$, one needs only specify the sets $Y_n$, and then show these sets contain the images of the restricted faces and degeneracies, since the definitions of the faces and degeneracies are forced by \cref{3.4.1}.
\end{remark}

\begin{definition}\label{3.4.4}
    Let $\alpha : \X\to\Y$ be a morphism of simplicial sets. The \emph{image} of $\alpha$, denoted $\im\alpha_\bullet$, is defined as the simplicial subset of $\Y$ given by
    \[
    (\im\alpha)_n=\im(\alpha_n : X_n\to Y_n),
    \]
    where $\alpha_n$ is the component of $\alpha$ at $n$. We make use of the convenient notation and write $\im\alpha_n$ to denote the set of $n$-simplices of $\im\alpha$.
\end{definition}

\begin{claim}\label{3.4.5}
    With notation as in \cref{3.4.4}, the image is indeed a simplicial subset of $\Y$.
\end{claim}
\begin{proof}
    By \cref{3.4.3}, we need only show the images of each face and degeneracy land within the specified sets. We start by showing this for degeneracies. Let $y\in\im\alpha_n$, and we want to show $s_i^{\im\alpha}(y)\in\im\alpha_{n+1}$. Recall by \cref{3.1.5} that morphisms of simplicial sets are natural tranformations, meaning they commute with degeneracies as indicated in the diagram
    \[\begin{tikzcd}
	{X_n} & {Y_n} \\
	{X_{n+1}} & {Y_{n+1}}
	\arrow["{\alpha_n}", from=1-1, to=1-2]
	\arrow["{s_i^X}"', from=1-1, to=2-1]
	\arrow["{s_i^Y}", from=1-2, to=2-2]
	\arrow["{\alpha_{n+1}}"', from=2-1, to=2-2]
    \end{tikzcd}\]
    By definition, for each $y\in\im\alpha_n$, there exists at least one $x\in X_n$ such that $y=\alpha_n(x)$. Since the degeneracy in $\im\alpha_\bullet$ is the same as the degeneracy for $\Y$, we have
    \[
    s^{\im\alpha}_i(y)=s^Y_i(y)=\left(s^Y_i\alpha_n\right)(x)=\left(\alpha_{n+1}s_i^X\right)(x),
    \]
    where the final step follows by commutativity of the diagram. By definition, this is in the image of $\alpha_{n+1}$, completing the proof for degeneracies. The proof for faces is the same, and is done by constructing the corresponding naturality diagram for faces instead of degeneracies.
\end{proof}

\begin{definition}\label{3.4.6}
    Let $\X,\Y$ be simplicial sets, and suppose that for each $n$, the faces and degeneraces of $\X$ and $\Y$ agree on the intersection $X_n\cap Y_n$. In symbols,
    \[
    s^Y_i|_{X_n\cap Y_n}=s_i^X|_{X_n\cap Y_n},\qquad d^Y_i|_{X_n\cap Y_n}=d_i^X|_{X_n\cap Y_n}.
    \]
    Then the \emph{union} $\X\cup\Y$ of $\X$ and $\Y$ is defined as the levelwise union
    \[
    (X\cup Y)_n=X_n\cup Y_n,
    \]
    and the faces and degeneracies are defined by
    \[
    s_i(x) = \begin{cases}
        s_i^X(x),&x\in X\\
        s_i^Y(x),&x\in Y
    \end{cases},\qquad d_i(x) = \begin{cases}
        d_i^X(x),&x\in X\\
        d_i^Y(x),&x\in Y
    \end{cases}.
    \]
\end{definition}

\begin{claim}\label{3.4.7}
    With notation as in \cref{3.4.6}, $\X\cup\Y$ is a simplicial set. In particular, $\X,\Y$ are simplicial subsets of their union.
\end{claim}
\begin{proof}
    The fact that $\X,\Y\subseteq\X\cup\Y$ follows manifestly from the definition, provided we indeed show $\X\cup\Y$ is a simplicial set, which is not terribly difficult. As usual, we appeal to \cref{3.1.7} and show that the faces and degeneracies are well defined, and satisfy the simplicial identities. Since $\X$ and $\Y$ agree on their intersection, for any $x\in X_n\cap Y_n$, we must have $d_i^X(x)=d_i^Y(x)$, and the same for degeneracies. Thus, the maps are well-defined. To show they satisfy the simplicial identities, note that for all $x\in X_n\cup Y_n$, we have $x\in X_n$ or $x\in Y_n$. Without loss of generality, let $x\in X_n$. Since $\X$ is a simplicial set, the simplicial operators satisfy the simplicial identities for $x$. Extending this argument to $\Y$, we find that the simplicial operators satisfy the simplicial identities for all $x\in X_n\cup Y_n$, and thus $\X\cup\Y$ is indeed a simplicial set.
\end{proof}

Previously in \cref{2.1.6}, we defined the \emph{face} of the standard $n$-simplex, and we were even able to encode this information into the category $\BD$ in \cref{2.4.4}. As explained in \cref{3.2.5}, the combinatorial and geometric information encoded by simplicial sets is in many ways more general than that of $\BD$, and as seen in \cref{3.2.7}, the standard $n$-simplicial sets encode the information of the standard geometric $n$-simplex. We should thus think that there should be some way to define the faces of the standard $n$-simplicial set. Indeed there is such a way to do so, and the definition of the face of the standard $n$-simplicial set will give rise to much more important constructions.

\begin{definition}\label{3.4.8}
    Let $d^i$ be the image of the coface map $[n-1]\to[n]$ under the Yoneda embedding. More precisely, let $d^i : \Delta^{n-1}\to\Delta^n$ denote precomposition by the coface map $d^i:[n-1]\to[n]$.
    \begin{enumerate}
        \item The $i$th face $\partial_i\Delta^n$ of the simplical set $\Delta^n$ is the simplicial subset $\im d^i_\bullet$.
        \item The \emph{boundary} $\partial\Delta^n$ of the standard $n$-simplicial set $\Delta^n$, also called the \emph{simplicial $n$-sphere}, is the simplicial subset
        \[
        \partial\Delta^n \coloneqq\bigcup_{i\in[n]} \partial_i\Delta^n=\bigcup_{i\in[n]}\im d^i_\bullet \subseteq\Delta^n.
        \]
        \item The \emph{$k$th horn} $\Lambda_n^k$ of the standard $n$-simplicial set is the simplicial subset
        \[
        \Lambda_n^k\coloneqq\bigcup_{i\in[n],i\neq k}\partial_i\Delta^n=\bigcup_{i\in[n],i\neq k}\im d^i_\bullet \subseteq\partial\Delta^n\subseteq\Delta^n.
        \]
    \end{enumerate}
\end{definition}

\begin{remark}\label{3.4.9}
    Because boundaries and horns are defined as unions over simplicial subsets of $\Delta^n$, their faces and degeneracies automatically agree with $\Delta^n$, and in particular agree with each other, making the union of simplicial sets well-defined.
\end{remark}

\begin{intuition}\label{3.4.10}
    The geometric intuition for faces has already been discussed. Our definition of faces implies that the coface maps $d^i : \Delta^{n-1}\to\Delta^n$ do the ``simplicial set equivalent'' of inserting the standard $(n-1)$-simplicial set at the ``$i$th face of the standard $n$-simplicial set''. It remains to be shown of course that the geometric realization of $\im d^i_\bullet$ is the $i$th face of $|\Delta^n|$, but this will be shown soon.

    For now, working with the understanding that $\im d^i_\bullet$ represents the face of the standard $n$-simplicial set, we define the \emph{boundary} of the standard $n$-simplicial set as the union of all of its faces. We expect a union of simplicial sets to correspond to a union of geometric objects, so working under this intuition, we would imagine the geometric realization of the boundary $\partial\Delta^n$ to be the physical boundary of the standard $n$-simplex. Once again, we will show this to be true shortly.

    The only definition that may not be immediately clear is that of a horn. The definition of the $k$th horn $\Lambda_n^k$ is the union of all the faces of $\Delta^n$ \emph{except} the $k$th one. For example, consider $\Lambda_2^2$. Geometrically, this corresponds to taking the standard $2$-simplex, deleting the interior so that we are only left with its faces (which constitute the boundary), and then deleting the $2$nd face. \Cref{fig:6} illustrates this geometrically.
    \begin{figure}[ht]
        \centering
        \begin{tikzpicture}[scale=5]
    % Set the viewing angle for the 3D plot
    \tdplotsetmaincoords{70}{110}

    % Establish the 3D coordinate system
    \begin{scope}[tdplot_main_coords]
    
        % Draw the axes for context
        \draw[-latex] (0,0,0) -- (1.2,0,0) node[anchor=north east]{$x$};
        \draw[-latex] (0,0,0) -- (0,1.2,0) node[anchor=north west]{$y$};
        \draw[-latex] (0,0,0) -- (0,0,1.2) node[anchor=south]{$z$};

        % Define the vertices of the standard 2-simplex
        \coordinate (v0) at (1,0,0);
        \coordinate (v1) at (0,1,0);
        \coordinate (v2) at (0,0,1);

        \draw[thick, blue] 
            (v1) -- (v2);
        \draw[thick, blue] 
            (v0) -- (v2);
        \draw[thin, dashed, gray] 
            (v0) -- (v1);

        % Label the vertices (0-faces)
        \node[at=(v0), anchor=south east, inner sep=2pt] {$v_0$};
        \node[at=(v1), anchor=south west, inner sep=2pt] {$v_1$};
        \node[at=(v2), anchor=south west, inner sep=5pt] {$v_2$};

        \node[at=($(v1)!0.5!(v2)$), anchor=south west, inner sep=3pt, text=blue] {$f_0$};
        \node[at=($(v0)!0.5!(v2)$), anchor=south east, inner sep=3pt, text=blue] {$f_1$};
        \node[at=($(v0)!0.5!(v1)$), anchor=north, inner sep=3pt, text=blue] {$f_2$ - Deleted face};

    \end{scope}
\end{tikzpicture}
        \caption{The second horn $|\Lambda_2^2|$ of $|\Delta^2|$}
        \label{fig:6}
    \end{figure}
    The remaining two faces, $f_0$ and $f_1$, assemble into a sort of ``$\wedge$'' shape, which looks somewhat like a two-dimensional horn. Hence the namesake. It remains to be shown of course that the geometric realizations align with these objects. We can in fact prove stronger results. In \cref{sec4.1}, we will show that geometric realization is a left adjoint functor, and thus preserves colimits. It thus suffices to show that unions and images can be viewed as colimits, and our intuition here will follow as \cref{3.4.13}.
\end{intuition}


\begin{proposition}\label{3.4.11}
    Let $F : \J\to\D^\C$ be a diagram with $\D$ complete, and let $\mathrm{ev}_c:\D^\C\to\D$ be the evaluation functor $F\mapsto F(c)$, $\alpha\mapsto\alpha_c$. Then $\lim F$, if it exists, is given by $(\lim F)(c)=\lim (\mathrm{ev}_c\circ F)$.
\end{proposition}
\begin{proof}
    See \cite[\S5.3~Thm.~1]{ml}.
\end{proof}

\begin{corollary}\label{3.4.12}
    Images and unions of simplicial sets, as defined in \cref{3.4.1} and \cref{3.4.6}, are colimits in $\sSet$. In particular, if $\X,\Y\in\sSet$ and $\alpha : \X\to\Y$ is a morphism of simplicial sets, then
    \begin{enumerate}
        \item $\left|\im(\alpha : \X\to\Y)_\bullet\right|\cong\im(\alpha : |\X|\to|\Y|)$, and
        \item $\left|\X\cup\Y\right|\cong|\X|\cup|\Y|$.
    \end{enumerate}
\end{corollary}
\begin{proof}
    Recall that $\sSet=\Set^{\BD\opp}$ is a functor category, and since $\Set$ is complete and cocomplete, \cref{3.4.11} and its dual apply. In particular, any diagram $F$ in $\sSet$ must satisfy $(\lim F)_n=\lim (\mathrm{ev}_n\circ F)$. This tells us that limit in $\sSet$ are computed levelwise, meaning in particular
    \begin{enumerate}
        \item Let $\alpha : \X\to\Y$ be a morphism of simplicial sets. The \emph{limit} definition of the image $\im\alpha_\bullet$ (not \cref{3.4.4}) must satisfy $(\im\alpha)_n=\im(\alpha_n)$, where the second is a limit in $\Set$. Note that this is precisely \cref{3.4.4}, and thus we have that our image in $\sSet$ can be regarded as a limit. Since $\Set$ admits canonical decompositions for morphisms, there is an isomorphism $\im(\alpha_n)\cong\coim(\alpha_n)$. The coimage is a colimit, and thus we may regard the image in $\sSet$ as isomorphic to a colimit.
        \item Let $\X,\Y$ be simplicial sets whose simplicial operators agree on their intersection. Because their simplicial operators agree, it is easy to see that $\X\cap\Y$ is a subsimplicial set of both $\X$ and $\Y$, with the intersection given by levelwise intersection. Note that the pushout $Z_\bullet$ of $\X$ and $\Y$ along the canonical inclusions $\X\cap\Y\hookrightarrow \X,\Y$ can be computed levelwise, as indicated in the diagram
        \[\begin{tikzcd}
	{(X\cap Y)_n} & {X_n} \\
	{Y_n} & {Z_n}
	\arrow[hook', from=1-1, to=1-2]
	\arrow[hook, from=1-1, to=2-1]
	\arrow["\lrcorner"{anchor=center, pos=0.125}, draw=none, from=1-1, to=2-2]
	\arrow[from=1-2, to=2-2]
	\arrow[from=2-1, to=2-2]
    \end{tikzcd}\]
    in $\Set$. It is well known that this pushout can be given by $Z_n=X_n\cup Y_n$, which is precisely the definition of the $n$-simplices of the union $(X\cup Y)_n$. Thus, we may regard the union as the pushout, which is in particular a colimit.
    \end{enumerate}
    Recall that geometric realiation $|-|$ is left adjoint, and thus preserves colimits. This implies that $|\X\cup\Y|\cong|\X|\cup|\Y|$ and that $|\coim\alpha|=\coim|\alpha|$. Note that $\Top$ also admits canonical decompositions of morphisms, so we obtain $\coim|\alpha|\cong\im|\alpha|$, giving a string of isomorphism
    \[
    |\im\alpha|\cong|\coim\alpha|\cong\coim|\alpha|\cong\im|\alpha|.
    \]
\end{proof}

\begin{corollary}\label{3.4.13}
    The following statements hold.
\begin{enumerate}
    \item The geometric realization of the $i$th face is the $i$th face of the geometric $n$-simplex:
    \[
    \left|\partial_i\Delta^n\right|\cong d^i(|\Delta^{n-1}|).
    \]
    \item The geometric realization of the boundary is the boundary of the geometric $n$-simplex:
    \[
    \left|\partial\Delta^n\right|\cong\partial\left|\Delta^n\right|.
    \]
    \item The geometric realization of the $k$th horn is the $k$th horn (defined as the union of all faces other than $k$) of the geometric $n$-simplex:
    \[
    \left|\Lambda_n^k\right|\cong\bigcup_{i\in[n],i\neq k}\im (d^i :\left|\Delta^{n-1}\right|\to\left|\Delta^n\right|).
    \]
\end{enumerate}
\end{corollary}
\begin{proof}
	Indeed, these are all special cases of $\cref{3.4.12}$, and it suffices to show $|d^i|=d^i : |\Delta^{n-1}|\to|\Delta^n|$. This is shown in \cref{sec3.5}.
\end{proof}

\begin{notation}\label{3.4.14}
    As is likely expected, we represent the geometric $k$th horn of the standard $n$-simplex as $|\Lambda_n^k|$. We also adopt the convention of denoting the $i$th face of the geometric $n$-simplex as $\partial_i|\Delta^n|$.
\end{notation}

\begin{remark}\label{3.4.15}
    It should be fairly clear to anyone who has taken topology that the boundary of the triangle $\partial|\Delta^2|$ is homeomorphic to the 1-sphere $S^1$. Similarly, it's clear that the boundary of the tetrahedron $\partial|\Delta^3|$ is homeomorphic to the 2-sphere $S^2$. Indeed this generalizes: for $n>1$, there is a homeomorphism $\partial|\Delta^n|\cong S^{n-1}$. For this reason, many authors call $\partial\Delta^n$ the simplicial $n$-sphere.
\end{remark}

We close this chapter by establishing a collection of important facts about horns. These facts are extremely useful in establishing basic results for \emph{Kan complexes} and \emph{quasi-categories}, which are explored a small amount in \cref{sec4.4}, and play a vital role in homotopy theory and higher category theory.

\begin{definition}\label{3.4.16}
	Let $X_{\bullet} $ be a simplicial set. A \emph{horn} of $X_{\bullet} $ is a morphism of simplicial sets $\Lambda_i^n \to X_{\bullet} $. A \emph{face} of $X_{\bullet} $ is a morphism of simplicial sets $\partial _i\Delta^n \to K$.
\end{definition}

\begin{lemma}\label{3.4.17}
	Let $X_{\bullet} $ be a simplicial set. Then a morphism $\sigma : \Delta^n \to X_{\bullet} $ of simplicial sets is simply a choice of an $n$-simplex $\tau\in X_n$, and mapping $\sigma(1_{[n]} )=\tau$.
\end{lemma}
\begin{proof}
	Indeed, this is precisely the Yoneda lemma applied to simplicial sets.
\end{proof}

\begin{remark}\label{3.4.18}
	Since the epi-mono factorization in $\boldsymbol{\Delta} $ given by \cref{2.3.4} is unique, and since elements of $\partial _i\Delta^n$ are simply given by postcomposing $\Delta^{n-1} $ by a coface, uniqueness gives an isomorphism $\partial_i\Delta^n\cong\Delta^{n-1} $ given by simply removing the coface from the composite. Applying \cref{3.4.17}, this gives us the fact that a morphism of simplicial sets $\partial _i\Delta^n \to X_{\bullet} $ is a choice of an $(n-1)$-simplex $\tau\in X_{n-1} $, and mapping $d^i \mapsto \tau$.
\end{remark}

\begin{terminology}
	In view of \cref{3.4.16,3.4.17,3.4.18}, we will often refer to morphisms $\Delta^n \to X_{\bullet} $ as $n$-simplices of $X_{\bullet} $, and will refer to morphisms $d^j\mapsto \tau\in X_n$ as \emph{faces} of $X_{\bullet} $.
\end{terminology}


\begin{proposition}\label{3.4.20}
	A horn $\Lambda_i^n \to X_{\bullet} $ in a simplicial set $X_{\bullet} $ is determined by an $(n-1)$-simplex $\tau_j$ for each $j\neq i,j\in[n]$ such that the compatibility condition $d_{j} (\tau_k)=d_{k-1} (\tau_j)$ for $j<k$ holds. The horn is then given by $d^j \mapsto \tau_j$.
\end{proposition}
\begin{proof}
	The proof is quite lengthy, so we offer only a sketch. In the proof of \cref{3.4.12}, we observed that unions in $\sSet $ are unions in the category-theoretic sense, so one constructs the relevant pullback diagram, which (by \cref{3.4.18}) simplifies the statement to ``the diagram commutes if and only if the compatibility condition holds''. If the diagram commutes, then the condition follows very easily by simplicial identities. For the converse direction, it suffices to show each face $d^j \mapsto \tau_j$ agrees on the intersection
	\[
		K_{\bullet} \coloneqq \bigcap_{j\in[n],j\neq i} \partial_j\Delta^n
	.\]
	This is shown by proving that $K_{\bullet} $ has $i$ as a unique vertex, and that it has $s_m s_{m-1} \ldots s_0(i)$ as a unique $m$-simplex. We then show the faces agree on $i$ by applying the compatibility condition on the $j$th and $(j+1)$th faces for $i<j$ and $j<i-1$; we then apply the condition for $i-1$ and $i+1$, showing any two faces agree by transitivity. By finiteness of $[n]$, induction gives that all faces agree on $i$, so it suffices to show they agree on degeneracies, which is immediate from the fact that the faces are morphisms of simplicial sets.
\end{proof}

\begin{definition}\label{3.4.21}
	Let $\sigma_0: \Lambda_i^n\to X_{\bullet} $ be a horn in $X_{\bullet} $. A \emph{filler} for the horn $\sigma_0$ is an $n$-simplex $\sigma : \Delta^n \to X_{\bullet} $ which agrees with $\sigma_0$ on the horn. More precisely,
	\[
		\sigma|_{\Lambda_i^n} =\sigma_0
	.\]
	This is often also expressed as the existence of the (not necessarily unique) dotted arrow as indicated
	\[\begin{tikzcd}[row sep = 2.5em, column sep = 2.5em]
	\Lambda_i^n \ar[" \sigma_0 ", r] \ar[hook, d]  & X_{\bullet}  \\
	\Delta^n \ar[" \sigma "', dotted, ur]  & 
	\end{tikzcd}\]
	rendering the diagram commutative.
\end{definition}


\begin{corollary}\label{3.4.22}
	Let $X_{\bullet} $ be a simplicial set. A filler for a horn $\sigma_0:\Lambda_i^n\to X_{\bullet} $ is an $n$-simplex $\sigma : \Delta^n \to X_{\bullet} $, $1_{[n]} \mapsto \tau\in X_n$, such that for each $j\neq i$ and each face $d^j : [n-1] \to [n]$,
			\[
				\sigma_0(d^j)=\sigma(d^j)
			.\]	
\end{corollary}
\begin{proof}
	The restriction $\sigma|_{\Lambda_i^n}$ gives a horn in $K$, and by \cref{3.4.20} this amounts to a choice $\tau_j$ for each $j\neq i$ satisfying the compatibiltiy condition, and mapping $d^j \mapsto \tau_j$. Also by \cref{3.4.20}, this horn agrees on (non-$i$) faces with $\sigma_0$ precisely when $\sigma|_{\Lambda_i^n} = \sigma_0$.
\end{proof}

\begin{exercise}[Challenge]\label{3.4.23}
	Technically, the reader now has the necessary knowledge to complete the following proof. This proof is fairly nontrivial, and the reader will likely struggle, especially if this is their first exposure to simplicial sets. We present the problem for the ambitious reader, but we do not expect the reader to solve it. A proof can be found in \cite[Prop.~1.1.2.2]{htt}. Let $X_{\bullet} $ be a simplicial set. Prove the following are equivalent:
	\begin{enumerate}
		\item There is a small category $\mathcal{C} $ and an isomorphism $N(\mathcal{C} )\cong X_{\bullet} $.
		\item Every \emph{inner horn} admits a \emph{unique} filler. More precisely, for all $0<i<n$, for all diagrams of solid arrows
			\[\begin{tikzcd}[row sep = 2.5em, column sep = 2.5em]
			\Lambda_i^n \ar[r] \ar[hook, d]  & X_{\bullet}  \\
			\Delta^n \ar[dashed, ur]  & 
			\end{tikzcd}\]
			there exists a \emph{unique} dashed arrow as indicated, rendering the diagram commutative.
	\end{enumerate}
\end{exercise}

